\documentclass[12pt,pdf,hyperref={unicode}]{beamer}
%\usetheme{boxes}
\beamertemplatenavigationsymbolsempty
\setbeamertemplate{footline}[page number]
% Set it for the internal PhD thesis defence to reduce number of slides
%\setbeamersize{text margin left=0.5em, text margin right=0.5em}

\usepackage[utf8]{inputenc}
\usepackage[T2A]{fontenc}
\usepackage[english, russian]{babel}
\usepackage{bm}
\usepackage{multirow}
\usepackage{ragged2e}
\usepackage{indentfirst}
\usepackage{multicol}
\usepackage{subfig}
\usepackage{amsmath,amssymb}
\usepackage{enumerate}
\usepackage{mathtools}
\usepackage{comment}
\usepackage[all]{xy}
\usepackage{tikz}
\usetikzlibrary{positioning,arrows}
\tikzstyle{name} = [parameters]
\definecolor{name}{rgb}{0.5,0.5,0.5}
\definecolor{beamerblue}{rgb}{0.20,0.20,0.75}
% ===== Пакет textpos подключён в преамбуле =====
\usepackage[absolute, overlay]{textpos}
\setlength{\TPHorizModule}{1mm}
\setlength{\TPVertModule}{1mm}

%\usepackage{caption}
%\captionsetup{skip=0pt,belowskip=0pt}

%\newtheorem{theorem}{Theorem}
%\newtheorem{statement}{Statement}
%\newtheorem{definition}{Definition}

% colors
\definecolor{darkgreen}{rgb}{0.0, 0.2, 0.13}
\definecolor{darkcyan}{rgb}{0.0, 0.55, 0.55}
\newcommand{\argmax}{\mathop{\mathrm{arg\,max}}}
%\AtBeginEnvironment{figure}{\setcounter{subfigure}{0}}
%\captionsetup[subfloat]{labelformat=empty}

% ------------------------------------------------------------
% Небольшие команды без добавления новых библиотек
% ------------------------------------------------------------
\newcommand{\argmin}{\mathop{\mathrm{arg\,min}}}
\newcommand{\R}{\mathbb{R}}
\newcommand{\HRule}{\rule{\linewidth}{0.3mm}}

\title{Детекция строк в фотографиях страниц рукописных текстов}
\author{Александр Забарянский}
\institute[]{Московский физико-технический институт}
\date{2026}

\begin{document}

% ------------------------------------------------------------
\begin{frame}
\titlepage
\end{frame}
\setcounter{page}{2}

% ------------------------------------------------------------
\begin{frame}{Цель исследования}

\vspace{3mm}

\begin{center}
\begin{minipage}{0.47\textwidth}
    \centering
    \includegraphics[width=\textwidth]{без_разметки.jpg}
    
    \vspace{1mm}
    \scriptsize Входное изображение
\end{minipage}
\hfill
\begin{minipage}{0.47\textwidth}
    \centering
    \includegraphics[width=\textwidth]{с_разметкой.jpg}
    
    \vspace{1mm}
    \scriptsize Результат детекции
\end{minipage}
\end{center}

% картинка на входе, картинка на выходе

Разработать алгоритм детекции строк в фотографиях страниц рукописных текстов,
устойчивый к \textbf{глобальному наклону}, \textbf{локальной кривизне} строки,
шуму бинаризации и неидеальной геометрии страницы.


\end{frame}

\begin{frame}{Постановка задачи}

\small

Набор изображений страниц с рукописным текстом:
\[
\mathcal{I}=\{I_k\}_{k=1}^K,
\qquad
I_k\in\mathbb{R}_+^{H_k\times W_k\times C}.
\]

Для каждого изображения задано множество строк:
\[
\mathcal{R}_k=\{R_{k,1},\dots,R_{k,n_k}\}.
\]

\vspace{1mm}
Найти отображение:
    \[
        f \colon \mathcal{I} \to 2^{B},
    \]
    \[
    B =
    \left\{
    \{p_i\}_{i=1}^{4}
    \;\middle|\;
    p_i=(x_i,y_i)\in \mathbb{Z}_{\ge 0}^{2},
    \text{ - вершинами прямоугольника}
    \right\}.
    \]

Необходимо, чтобы:
    \[
        \forall  R = \{p_i\}_{i=1}^{4} \in \hat{\mathcal{R}},
        \qquad p_i=(x_i,y_i),
    \]
    выполнялось: %условие принадлежности всех его вершин области изображения:
    \[
        p_i \in \Omega_I,
        \qquad
        \Omega_I = \{0,\dots,W\}\times\{0,\dots,H\},
        \qquad i=1,\dots,4.
    \]
    где $H$ и $W$ — высота и ширина изображения.

\end{frame}
% ------------------------------------------------------------
\begin{frame}{Датасет и синтетические данные}

\small

\begin{block}{Основной датасет}
Используется датасет HWR200: фотографии страниц с рукописным русским текстом. С помощью модели PaddleOCR-VL (1.5)
проведена разметка.
\end{block}

\vspace{1mm}

\begin{block}{Синтетика}
Дополнительно генерируем синтетические страницы:
\begin{itemize}
    \item $16000$ страниц размера $1400 \times 1700$;
    \item текст берётся из русских слов и рендерится handwritten-fonts;
    \item на странице от $10$ до $64$ строк;
    \item добавляется глобальная деформация страницы:
    парабола или загиб одного края.
\end{itemize}
\end{block}

\let\thefootnote\relax
\footnotetext{\href{https://huggingface.co/datasets/AntiplagiatCompany/HWR200}{HWR200: Russian handwritten text dataset.}}

\end{frame}


% ------------------------------------------------------------
\begin{frame}{Оценка качества алгоритма}

\footnotesize
%Intersection over Union}
Для каждой пары предсказанного и истинного полигонов:
\[
\operatorname{IoU}(\hat{R}_i,R_j)
=
\frac{|\hat{R}_i\cap R_j|}
     {|\hat{R}_i\cup R_j|}.
\]


%Для каждого предсказанного полигона $\hat{R}_i$ выбирается ещё не занятый
%ground truth полигон с максимальным IoU:
\[
j^*=\argmax_{j\in\mathcal{R}_{\mathrm{free}}}
\operatorname{IoU}(\hat{R}_i,R_j).
\]

Если
\[
\operatorname{IoU}(\hat{R}_i,R_{j^*})\ge 0.5,
\]
то пара считается $TP$. Иначе предсказание считается $FP$.
Все несопоставленные полигоны считаются $FN$.

\vspace{1mm}

\[
\operatorname{precision}=\frac{TP}{TP+FP},
\qquad
\operatorname{recall}=\frac{TP}{TP+FN},
\]
\[
\operatorname{hmean}
=
\frac{
2\cdot\operatorname{precision}\cdot\operatorname{recall}
}{
\operatorname{precision}+\operatorname{recall}
}.
\]

\end{frame}

% ------------------------------------------------------------
\begin{frame}{Метод 1: метод Хафа}

\small

\vspace{1mm}

\begin{enumerate}
    \item Бинаризуем изображение и выделяем связные компоненты.
    \item Оцениваем среднюю высоту символа $AH$ и ширину $AW$.
    \item Делим компоненты на три множества:
    \[
    \mathcal{C}_1 = \{c: 0.5AH \le h_c < 3AH,\; w_c \ge 0.5AW\},
    \]
    \[
    \mathcal{C}_2 = \{c: h_c \ge 3AH\},
    \qquad
    \mathcal{C}_3 = \mathcal{C}\setminus(\mathcal{C}_1\cup\mathcal{C}_2).
    \]
    \item Для компонент из $\mathcal{C}_1$ строятся block-centers.
    \item Эти точки голосуют в пространстве Хафа:
    \[
        \rho = x\cos\theta + y\sin\theta.
    \]
    \item Найденные линии используются для группировки компонент в строки.
\end{enumerate}

\let\thefootnote\relax
\footnotetext{\href{https://doi.org/10.1016/j.patcog.2008.05.011}{Louloudis et al., Text line detection in handwritten documents, Pattern Recognition, 2008.}}

\end{frame}


% ------------------------------------------------------------
\begin{frame}{Метод 1: наши изменения метода Хафа}

\small

\begin{block}{Что меняем}
\begin{enumerate}
    \item При оценке $AH$ отсекаем первый bin гистограммы высот.
    \item До голосования Хафа разделяем компоненты из \texttt{subset\_2}.
    \item Диапазон углов Хафа задаём относительно глобального угла, найденного через HPP:
    \[
    \theta^*=\argmax_{\theta\in[-30^\circ,30^\circ]}
    \operatorname{Var}\left(\sum_x B_\theta(x,y)\right).
    \]
    \item Убираем стадию создания новых строк.
\end{enumerate}
\end{block}

\vspace{1mm}

\begin{block}{Результат}
\centering
\scriptsize
\begin{tabular}{|l|c|c|c|}
\hline
Метод & Precision & Recall & Hmean \\
\hline
Paper method & 0.4519 & 0.5472 & 0.4950 \\
Our method   & 0.7426 & 0.7955 & 0.7681 \\
\hline
\end{tabular}
\end{block}

\end{frame}


% ------------------------------------------------------------
\begin{frame}{Метод 2: HPP + seam carving}

\small

\vspace{1mm}

\begin{enumerate}
    \item По бинарному изображению строится HPP:
    \[
        H(y)=\sum_{x=0}^{W-1} B(x,y).
    \]
    \item Профиль нормализуется:
    \[
        H_{\text{norm}}(y)=
        \frac{H(y)-\mu_H}{\sigma_H}.
    \]
    \item Строки выбираются по порогу:
    \[
        H_{\text{norm}}(y) > \tau.
    \]
    \item Соседние строки объединяются в регионы.
    \item Между соседними регионами запускается поиск минимального разреза.
\end{enumerate}

\let\thefootnote\relax
\footnotetext{\href{https://www.sciencedirect.com/science/article/pii/S2590123023002372}{Das and Panda, Seam carving, HPP and contour tracing, Results in Engineering, 2023.}}

\vspace{1mm}

%Метод почти не учитывает наклон строк и слишком грубо анализирует HPP.

\end{frame}


% ------------------------------------------------------------
\begin{frame}{Метод 2: наши изменения HPP}

\small

\begin{block}{Что меняем}
\begin{enumerate}
    \item Сначала исправляем глобальный наклон:
    \[
    \theta^*=\argmax_{\theta\in[-30^\circ,30^\circ]}
    \operatorname{Var}\left(\sum_x B_\theta(x,y)\right).
    \]
    \item Затем делаем локальное выпрямление.
    \item HPP анализируем после морфологического сглаживания:
    $
        H_{\text{smoothed}}
        =
        \operatorname{open}(\operatorname{close}(\tilde H,k),k).
    $
    \item Проверяем два ядра:
    $
        k_1=\frac{AH}{2}, k_2=\frac{AH}{3}.
    $
    %и выбираем вариант с большим числом найденных строк.
\end{enumerate}
\end{block}

\vspace{1mm}

\begin{block}{Результат}
\centering
\scriptsize
\begin{tabular}{|l|c|c|c|}
\hline
Метод & Hmean & Precision & Recall \\
\hline
Paper method & 0.7337 & 0.8492 & 0.6458 \\
Our method   & 0.8419 & 0.8494 & 0.8345 \\
\hline
\end{tabular}
\end{block}

\end{frame}

\begin{frame}{Локальное выпрямление кривизны}

\small

\begin{textblock*}{72mm}(6mm,16mm)
1. Оценка локального угла.

Для каждой ячейки находим
$\alpha(x,y)$.

\vspace{1.5mm}
2. Оценка локального наклона.

\hspace*{2mm}$s(x,y)=\tan\alpha(x,y)$

\vspace{1.5mm}
3. Интегрирование наклона по горизонтали.

\hspace*{2mm}$d_y(x,y)=\sum_{t=0}^{x-1}s(t,y)$

\vspace{1.5mm}
4. Центрирование смещений по каждой строке.

\hspace*{2mm}$d_y(x,y)\leftarrow d_y(x,y)-\frac{1}{W}\sum_{x=0}^{W-1}d_y(x,y)$

\vspace{1.5mm}
5. Построение remap-преобразования.

\hspace*{2mm}$T(x,y)=(x,\; y+d_y(x,y))$

\vspace{1.5mm}
6. Получение выпрямленного изображения.

\hspace*{2mm}$B_{\text{warp}}(x,y)=B\!\left(T(x,y)\right)$

\vspace{2mm}
\textbf{Идея.}
Локальные углы интерполируются на всё изображение,
после чего по карте вертикальных смещений выполняется
геометрическое выпрямление бинарной страницы.
\end{textblock*}

\begin{textblock*}{40mm}(82mm,14mm)
\centering
\includegraphics[width=40mm,height=28mm,keepaspectratio]{карта_углов.jpg}
\end{textblock*}

\begin{textblock*}{40mm}(82mm,43mm)
\centering
\includegraphics[width=40mm,height=28mm,keepaspectratio]{карта_смещений.jpg}
\end{textblock*}

\begin{textblock*}{40mm}(82mm,72mm)
\centering
\includegraphics[width=40mm,height=28mm,keepaspectratio]{итоговое_выпрямление.jpg}
\end{textblock*}

\end{frame}


% ------------------------------------------------------------
\begin{frame}{Метод 3: DBNet++}

\small

\begin{center}
    \includegraphics[width=0.75\textwidth]{Dbnet.jpg}
\end{center}

\begin{enumerate}
    \item Сеть строит вероятностную карту текста:
    \[
        P(x,y)\in[0,1].
    \]
    \item Дополнительно предсказывается адаптивный порог:
    \[
        T(x,y)\in[0,1].
    \]
    \item Дифференцируемая бинаризация:
    \[
        \hat{B}(x,y)=
        \frac{1}{1+\exp(-k(P(x,y)-T(x,y)))}.
    \]
    \item По бинарной карте строятся итоговые полигоны строк.
\end{enumerate}


\end{frame}


% ------------------------------------------------------------
\begin{frame}{Метод 3: DBNet++ и синтетика}

\small

\begin{block}{Что сравниваем}
Сравниваем DBNet++ без синтетики и DBNet++ с добавлением синтетических страниц.
\end{block}

\vspace{2mm}

\begin{block}{Результат}
\centering
\scriptsize
\begin{tabular}{|l|c|c|c|}
\hline
Метод & Hmean & Precision & Recall \\
\hline
DBNet++ без синтетики & 0.9567 & 0.9817 & 0.9329 \\
DBNet++ с синтетикой  & 0.9577 & 0.9805 & 0.9359 \\
\hline
\end{tabular}
\end{block}

\vspace{1mm}

\begin{block}{Вывод}
Синтетика немного увеличивает recall и итоговый Hmean,
то есть модель начинает пропускать меньше строк.
\end{block}

\let\thefootnote\relax
\footnotetext{\href{https://arxiv.org/abs/2202.10304}{Liao et al., Real-Time Scene Text Detection with Differentiable Binarization and Adaptive Scale Fusion, 2022.}}

\end{frame}


% ------------------------------------------------------------
\begin{frame}{Выводы и дальнейшая работа}

\small

\begin{block}{Выводы}
\begin{enumerate}
    \item Для метода Хафа и HPP удалось заметно улучшить precision, recall и Hmean.
    \item Лучшее качество показал DBNet++:
    \[
        Hmean = 0.9577.
    \]
\end{enumerate}
\end{block}

\vspace{1mm}

\begin{block}{Дальнейшая работа}
\begin{enumerate}
    \item Сравнить DBNet++ с другими современными нейросетевыми детекторами строк.
    \item Оптимизировать предварительную обработку для улучшения методов обучения без учителя.
\end{enumerate}
\end{block}

\end{frame}

% ------------------------------------------------------------
\begin{frame}{Источники}

\footnotesize

\begin{enumerate}
    \item Hough-based text line extraction: \\
    \href{https://ieeexplore.ieee.org/document/602034}
    {Likforman-Sulem et al., ICDAR 1995}

    \item Text line detection in handwritten documents: \\
    \href{https://doi.org/10.1016/j.patcog.2008.05.011}
    {Louloudis et al., Pattern Recognition, 2008}

    \item HPP + seam carving: \\
    \href{https://www.sciencedirect.com/science/article/pii/S2590123023002372}
    {Das and Panda, Results in Engineering, 2023}

    \item DBNet++: \\
    \href{https://arxiv.org/abs/2202.10304}
    {Liao et al., Differentiable Binarization and Adaptive Scale Fusion, 2022}

    \item HWR200 dataset: \\
    \href{https://huggingface.co/datasets/AntiplagiatCompany/HWR200}
    {HWR200: Russian handwritten text dataset}
\end{enumerate}

\end{frame}


% \begin{frame}{Постановка задачи} \begin{block}{Дано} $\mathcal{I}=\{I_k\}_{k=1}^K$, --- множество изображений рукописного текста, где $I_k \in \mathbb{R}_+^{H_k \times W_k \times C}$ Для каждого изображения $I_k$ задано множество $\mathcal{R}_k=\{R_{k,1},\dots,R_{k,n_k}\}$, где $R_{k,j}$ --- минимальный ограничивающий прямоугольник $j$-й строки. \end{block} \vspace{0.5mm} \begin{block}{Найти} Построить отображение $f:\mathcal{I}\to 2^B$, где $B=\{(x_1,y_1),(x_2,y_2),(x_3,y_3),(x_4,y_4)\mid x_i,y_i\in\mathbb{Z}_{\ge0}\}$. \end{block} \vspace{0.5mm} \begin{block}{Критерий} $\mathrm{Score}= \frac{1}{K}\sum_{k=1}^K \frac{1}{\max(n_k,\hat n_k)} \sum_{(i,j)\in \mathrm{matching}_k}\mathrm{IoU}(i,j)$, где $\mathrm{matching}_k$ --- оптимальное соответствие для $k$-го изображения. \end{block} \end{frame} 
%
% \begin{frame}{Решение}
%
% \begin{tikzpicture}[remember picture, overlay]
%
%     \def\imgw{36mm}
%     \def\topY{-11mm}
%     \def\midArrowY{-33mm}
%     \def\bottomY{39mm}
%     \def\bottomArrowY{27mm}
%     \def\lineY{15mm}
%
%     \node[anchor=north west, inner sep=0pt] 
%         (top1) at ([xshift=4mm,yshift=\topY]current page.north west)
%         {\includegraphics[width=\imgw]{input_image.jpg}};
%
%     \node[anchor=north, inner sep=0pt]
%         (top2) at ([yshift=\topY]current page.north)
%         {\includegraphics[width=\imgw]{binary.jpg}};
%
%     \node[anchor=north, inner sep=0pt]
%         at ([xshift=-20mm, yshift=\midArrowY]current page.north)
%         {\includegraphics[width=9mm]{чёрная_стрелка.jpg}};
%
%     \node[anchor=south west, inner sep=0pt]
%         (mid1) at ([xshift=4mm,yshift=\bottomY]current page.south west)
%         {\includegraphics[width=\imgw]{angle_map.jpg}};
%
%     \node[anchor=south, inner sep=0pt]
%         (mid2) at ([yshift=\bottomY]current page.south)
%         {\includegraphics[width=\imgw]{correct_binary.jpg}};
%
%     \node[anchor=south, inner sep=0pt]
%         at ([xshift=-20mm, yshift=29mm]current page.south)
%         {\includegraphics[width=9mm]{чёрная_стрелка.jpg}};
%
%     \node[anchor=south west, inner sep=0pt]
%         (line1) at ([xshift=4mm,yshift=\lineY]current page.south west)
%         {\includegraphics[width=\imgw]{seems_lines.jpg}};
%
%     \node[anchor=south, inner sep=0pt]
%         (line2) at ([yshift=\lineY]current page.south)
%         {\includegraphics[width=\imgw]{segment_lines.jpg}};
%
% \end{tikzpicture}
%
% \let\thefootnote\relax
% \footnotetext{\href{https://www.sciencedirect.com/science/article/pii/S2590123023002372?via\%3Dihub}{Seam carving, horizontal projection profile.}}
% \footnotetext{\href{https://www.mdpi.com/2313-433X/10/3/65}{Historical text line segmentation using U-Net networks.}}
%
% \end{frame}
% % ------------------------------------------------------------
% \begin{frame}{Вычислительный эксперимент}
%
% \begin{columns}[T]
% \begin{column}{0.46\textwidth}
%
%
% \begin{block}{Коррекция строк}
% Для изображения оценивается единый угол поворота:
% \[
% \theta^*=\argmax_{\theta\in[\theta_{\min},\,\theta_{\max}]}
% \mathrm{Var}\bigl(h_\theta(y)\bigr).
% \]
% \end{block}
%
% \vspace{1mm}
%
% \end{column}
%
% \begin{column}{0.54\textwidth}
%
% \centering
% \includegraphics[width=0.82\linewidth,height=28mm,keepaspectratio]{прибер_без_геометрической корекции.jpg}
%
%
% \vspace{3mm}
%
% \includegraphics[width=0.82\linewidth,height=28mm,keepaspectratio]{пример_улучшения.jpg}
%
%
% \end{column}
% \end{columns}
%
% \vspace{2mm}
%
% \begin{block}{Результаты}
% \centering
% \scriptsize
% \begin{tabular}{|l|c|c|}
% \hline
% Метрика & Коррекция строк & Геом. выпрямление \\
% \hline
% correctly detected lines & 31.67 & 33.67 \\
% \hline
% \end{tabular}
% \end{block}
%
% \end{frame}
%
% % ------------------------------------------------------------
% \begin{frame}{Заключение}
%
% \begin{enumerate}
%     \item Реализован алгоритм детекции строк в фотографиях рукописных страниц.
%     \item Предложено \textbf{локальное выпрямление страницы}.
%     \item Проведён эксперимент, демонстрирующий преимущество предложенного метода выпрямления по сравнению с простой коррекцией перспективы.
% \end{enumerate}
%
% \end{frame}

\end{document}